
\begin{saveblock}{codebox}%
	\begin{minted}{bib}
		author = {A. Smith and B. Doe and E. Dropper and F. Laker}
	\end{minted}
\end{saveblock}



\begin{frame}{Bibfile-meerdere auteurs}
	
	In je \mintinline{tex}|.bib|-bestand, scheid auteurs met \mintinline{tex}|and|:
	\useblock{codebox}
	Zo kan \mintinline{tex}|biblatex| controleren hoeveel auteurs het toont.\\ \pause
	Dit is belangrijk voor het command \mintinline{tex}|\textcite{source}|:
	\begin{itemize}[label=\textbullet]
		\item \mintinline{tex}|\usepackage{biblatex}| geeft Smith et all \pause
		\item \mintinline{tex}|\usepackage[maxcitenames=5]{biblatex}| geeft Smith, Doe, Dropper and Laker \pause
		\item \mintinline{tex}|\usepackage[maxcitenames=3, mincitenames=2]{biblatex}| geeft Smith, Doe et all \pause
	\end{itemize}
	maxbibnames en minbibnames doen hetzelfde voor je literatuurlijst


	%Voorbeeld: Peter Adams et al.	
\end{frame}

\begin{saveblock}{codebox}%
	\begin{minted}{bib}
		@article{Smith1994,
			author = {A. Smith and B. Doe and E. Dropper},
			title = {Titel},
			journal = {Opt. Lett.},
			year = 1994,
			volume = 19,
			number = 11,
			pages = {780-782}
		}
	\end{minted}
\end{saveblock}


\begin{saveblock}{codeboxB}%
	\begin{minted}{bib}
		@article{Smith1994,
			author = {A. Smith and B. Doe and E. Dropper},
			title = {Titel},
			journal = {Opt. Lett.},
			year = 1994,
			volume = 19,
			number = 11,
			pages = {780-782}
		}
	\end{minted}
\end{saveblock}

